% ACL long paper. Numbers: analysis/acl_rewrite_stats.py (OfficeBench, computed from run logs)
% and the four failure-analysis reports (tau2-bench, WorkBench, TheAgentCompany).
\documentclass[11pt]{article}
\usepackage[review]{acl}
\usepackage{times}
\usepackage{latexsym}
\usepackage[T1]{fontenc}
\usepackage[utf8]{inputenc}
\usepackage{microtype}
\usepackage{graphicx}
\usepackage{booktabs}
\usepackage{array}
\usepackage{amsmath}
\usepackage{enumitem}
\usepackage{xcolor}
\usepackage{url}

% ---- verification notes: set \showverifyfalse to hide every note ----
\newif\ifshowverify
\showverifytrue
\newcommand{\verify}[1]{\ifshowverify{\color{red}\,[VERIFY: #1]}\fi}

\newcolumntype{L}[1]{>{\raggedright\arraybackslash}p{#1}}
\newcommand{\gemma}{\texttt{gemma4-31b}}
\newcommand{\gptoss}{\texttt{gpt-oss-20b}}
\newcommand{\OB}{OfficeBench}
\newcommand{\WB}{WorkBench}
\newcommand{\TAU}{$\tau^2$-bench}
\newcommand{\TAC}{TheAgentCompany}

\title{Is It the Agent or the Pipeline?\\
Attributing English--Bangla Failures of LLM Agents across Four Benchmarks}

\author{Anonymous ACL submission}

\begin{document}
\maketitle

\begin{abstract}
LLM agents are evaluated mostly in English. When they are evaluated in another language, a lower
score is usually read as a weaker agent. But the score is produced by a pipeline: a translated
task, sometimes a simulated user, tools and an automatic grader. We run two open-weight agents,
\gemma{} and \gptoss{}, on paired English and Bangla versions of four benchmarks
(\OB{}, \WB{}, \TAU{} and \TAC{}) and trace every task that passes in English but fails in
Bangla. The Bangla gap depends on model and benchmark: after Holm correction, \gemma{} loses
8.0 points on \OB{} and 11.5 on \TAU{}; no other comparison is significant. Of 163
English-pass/Bangla-fail cases, the largest group, 74 (45.4\%), traces to evaluation components:
a Bangla-speaking simulated user, English-only checkers, translation or the environment. Between
14.7\% and 54.6\% of cases can be Bangla-specific agent failures, depending on how unclear and
shared failures are counted. Crediting answers rejected only for their script cuts \gemma{}'s
\OB{} gap from 8.0 to 3.0 points. The genuine failures on \OB{} share one pattern: the agent
keeps the meaning but writes the wrong form. Raw cross-lingual agent scores can overstate
language-specific deficits.
\end{abstract}

\section{Introduction}

\begin{figure*}[t]
\centering
\includegraphics[width=\textwidth]{figures/fig_overview.pdf}
\caption{Evaluation protocol. (a) Each task runs once in English and once in Bangla; only
natural-language text changes. (b) The agent loop. (c) The benchmark's grader scores each run, and
the two runs of a task form one paired outcome. (d) Every English-pass/Bangla-fail case is traced
and assigned the first cause whose check fails (codes in Table~\ref{tab:attribution} and
Appendix~\ref{app:codebook}).}
\label{fig:overview}
\end{figure*}

LLM agents plan over many steps, call tools and change files, databases and conversations
\citep{yao2023react,liu2023agentbench,zhou2024webarena}. Their benchmarks are almost all in
English \citep{jimenez2024swebench,xie2024osworld,mialon2024gaia}. Recent work translates agent
benchmarks into other languages and reports lower scores
\citep{hofman2026maps,wang2025xwebagentbench,li2026polyworkbench}. The usual reading is that the
agent is weaker in that language.

That reading assumes the score measures only the agent. It does not. A score is produced by a
pipeline: a translated task, sometimes a simulated user or simulated coworkers, tools, and an
automatic grader. Any of these can fail in the new language. Two examples from our runs show
this. On \OB{}, \gemma{} answers a Bangla question with the correct number written in Bengali
digits; the checker searches for the ASCII string ``40'' and marks the run as failed. On \TAU{},
the agent asks the Bangla-speaking simulated customer to confirm a cancellation; the customer
confirms and ends the conversation in the same message, so the agent never performs the
cancellation. Both runs count as Bangla failures. Neither is a failure of the agent.

We therefore ask where English--Bangla failures come from. We run two open-weight agents,
\gemma{} and \gptoss{}, on paired English and Bangla versions of four benchmarks that route the
outcome through different components (Figure~\ref{fig:overview}). For every task that passes in
English but fails in Bangla, we compare the two trajectories and assign the failure to one cause.
We address three questions:

\begin{itemize}[noitemsep,topsep=2pt,leftmargin=*]
\item \textbf{RQ1.} How large is the English--Bangla gap, and does it depend on the model and the
benchmark?
\item \textbf{RQ2.} Where in the trajectory do Bangla-only failures first occur?
\item \textbf{RQ3.} Which pipeline component is responsible for each Bangla-only failure?
\end{itemize}

\noindent Our contributions are:
\begin{itemize}[noitemsep,topsep=2pt,leftmargin=*]
\item A paired English--Bangla evaluation of two agents on four benchmarks, with discordant
counts, exact tests and bootstrap intervals for every comparison (\S\ref{sec:rq1}).
\item A case-level attribution protocol with operational definitions and a fixed decision order,
applied to all 163 English-pass/Bangla-fail cases (\S\ref{sec:attribution}, \S\ref{sec:rq3}).
\item Evidence that the Bangla gap is model- and benchmark-specific, that the largest group of
Bangla-only failures traces to evaluation components, and that correcting the grading removes
most of one significant gap (\S\ref{sec:rq3}).
\item Concrete recommendations for builders of multilingual agent benchmarks (\S\ref{sec:discussion}).
\end{itemize}

\section{Related Work}

\paragraph{Agent benchmarks.} Interactive benchmarks evaluate LLMs as agents
\citep{liu2023agentbench,zhou2024webarena,xie2024osworld,jimenez2024swebench,mialon2024gaia}.
We use four with different outcome paths: \OB{} grades files produced across office applications
\citep{wang2024officebench}; \WB{} grades the final database after a single request
\citep{styles2024workbench}; \TAU{} adds an LLM-simulated user \citep{barres2025tau2}, extending
$\tau$-bench \citep{yao2024taubench}; and \TAC{} runs long company tasks with simulated coworkers
in an OpenHands harness \citep{xu2024theagentcompany,wang2025openhands}.

\paragraph{Multilingual evaluation of LLMs and agents.} Static multilingual benchmarks show an
English advantage that grows for lower-resource languages
\citep{conneau2018xnli,ahuja2023mega,shi2023mgsm,bandarkar2024belebele,huang2023notall,zhang2023dont,etxaniz2024multilingual}.
Recent agent benchmarks translate tasks and measure the gap: MAPS covers four agent suites in
eleven languages \citep{hofman2026maps}, and X-WebAgentBench, OmnilingualGAIA2, PolyWorkBench,
SEATauBench and MASSIVE-Agents cover web, assistant, workplace, conversational and
function-calling tasks
\citep{wang2025xwebagentbench,caciolai2026omnilingualgaia2,li2026polyworkbench,nguyen2026seataubench,kulkarni2025massiveagents}.
These works quantify the size of the gap. We ask which component of the pipeline produces it.

\paragraph{Multilingual tool use.} Agents must emit well-formed calls
\citep{schick2023toolformer,qin2024toolllm,patil2024gorilla}. In multilingual settings, models
often write argument values in the user's language instead of the form the tool expects
\citep{luo2026lost,chauhan2026api}. These studies isolate argument-language errors in tool
calling. We find a related form error inside full agent runs, and we separate it from errors of
the checker and the simulated user.

\paragraph{Validity and failure attribution.} The Agentic Benchmark Checklist
\citep{zhu2025abc} and protocol-validity audits \citep{shao2026protocol} show that setup and
grading flaws can change conclusions. Failure taxonomies \citep{cemri2025mast} and diagnostic
tools \citep{barke2026agentrx} analyse trajectories, and automatic attribution of the decisive
step remains unreliable (53.5\% agent-level and 14.2\% step-level accuracy;
\citealp{zhang2025whichagent}). We apply manual, evidence-based attribution to cross-lingual
failures and rule out evaluation causes before attributing a failure to the agent.

\section{Methodology}
\label{sec:method}

\subsection{Benchmarks and sampling}
Table~\ref{tab:benchmarks} summarises the four benchmarks. We use \OB{} in full (300 tasks).
For \WB{} we draw 15 tasks from each of its six domains with a fixed seed (90 of 690 tasks)
\verify{seed value}. For \TAU{} we select tasks from each domain in their original order and drop
every task whose English or Bangla run ended in an infrastructure or quota error. This leaves 218
paired tasks for \gemma{} and 149 for \gptoss{}; 149 tasks are gradeable in all four cells
\verify{the draft's Table~1 gives 149 tasks for \TAU{}; the number of tasks sampled before
exclusions is not recorded in the sources}. For \TAC{} we run 23 English--Bangla task pairs over 21
unique tasks; two tasks appear in both of our two task sets, and two further sampled tasks ran only
in English and are excluded.

\begin{table*}[t]
\centering\footnotesize
\setlength{\tabcolsep}{2.5pt}
\begin{tabular}{L{2.1cm}L{3.2cm}L{3.2cm}L{3.2cm}L{3.2cm}}
\toprule
 & \OB{} & \WB{} & \TAU{} & \TAC{} \\
\midrule
Interaction & multi-step & single-turn & dialogue & long-horizon \\
Simulated & -- & -- & customer & coworkers \\
Grading & file checks & DB state & DB state, actions, device & checkpoints \\
Sample & full & 15 per domain & ordered, per domain & subset \\
Paired $n$ & 300 & 90 & 218 / 149 & 23 (21 tasks) \\
Models & both & both & both & \gemma{} \\
\midrule
In Bangla & system prompt, task & system prompt, task, reply & scenario, policy, agent prompt & task, coworker replies \\
Kept English & tools, IDs, files, graders & tools, IDs, dates & IDs, product names, enums & apps, files, URLs \\
\bottomrule
\end{tabular}
\caption{Benchmarks. Paired $n$: tasks run and graded in both languages (\TAU{}: \gemma{} /
\gptoss{}). The localization rows describe what was translated in each benchmark
\verify{the draft states that policies stay in English; the \TAU{} report lists the policy as
translated}.}
\label{tab:benchmarks}
\end{table*}

\subsection{Bangla localization}
\label{sec:localization}
In every benchmark the natural-language task is in Bangla, and everything the environment
executes stays in English: tool names, argument keys, identifiers and the grader
(Table~\ref{tab:benchmarks}). This keeps every action executable. It does not make language the
only difference between the two runs. The translation itself can change a task, and on two
benchmarks a simulated user or coworkers speak Bangla. We treat both as part of the Bangla
condition and measure their effect in \S\ref{sec:rq3}.

Translations were LLM-assisted. The intended rules were: keep names, URLs and file paths in Latin
script; keep strings matched by a grader in English; write all other content, including numbers
and dates, as a Bangla speaker would \verify{the draft states that an independent review checked
meaning; reviewer and coverage are not recorded}. Writing numbers in Bangla style and letting the
agent answer in Bangla are part of the treatment, and both can cause English graders to reject
correct work. For \OB{}, the instructions were machine-translated by \gemma{} and cached, so both
agents saw the same string. An audit of the 298 translated instructions shows that the rules were
applied unevenly: person names stayed in Latin script in 43 of 103 mentions, 44 instructions use
Bengali digits and 50 use ASCII digits (14 use both), and 1 of 262 file names was altered.

\subsection{Experimental setup}
Both models were served through an Ollama endpoint with temperature 0 (Appendix~\ref{app:repro}).
Each benchmark runs with its own harness and limits: \OB{} allows 20 steps per task, \WB{} uses
single-turn function calling, \TAU{} pairs the agent with a simulated customer played by the same
model, and \TAC{} runs an OpenHands agent \verify{100-iteration cap stated in the draft; not in
the sources}. Each task is run once per language with identical settings, so every comparison is
paired at the task level; on \TAC{}, 13 Bangla tasks were also rerun, and we use the latest run.
We log every instruction, model output, tool call, tool result, error and final answer.

For \OB{} we repaired three defects in the shipped scorer: email checks now read the run's own
files, mailbox lookup tries the exact account name first, and text files are read as UTF-8. No
check was relaxed. The repairs raise the pass counts from 132, 103, 80 and 97 to 151, 127, 86 and
106 (\gemma{} English and Bangla, \gptoss{} English and Bangla) and turn no pass into a fail.

\subsection{Measuring the gap}
Let EN and BN be the success rates in English and Bangla. We report the absolute gap
$\Delta = \mathrm{EN} - \mathrm{BN}$ and the multilingual effectiveness loss rate,
$\mathrm{MELR} = (\mathrm{EN}-\mathrm{BN})/\mathrm{EN}$, the share of English successes lost in
Bangla. MELR is unstable when EN or the number of differing tasks is small, so we always report it
with $n$. A \emph{discordant pair} is a task with different outcomes in the two languages: an
\emph{EN-pass/BN-fail} (EN-only) or a \emph{BN-pass/EN-fail} (BN-only) task. We test each gap
with McNemar's exact test on the discordant pairs \citep{mcnemar1947}, apply Holm correction
across the seven model--benchmark comparisons \citep{holm1979}, and report 95\% paired bootstrap
intervals (4{,}000 resamples of tasks). Subgroup tests are Holm-corrected within their family.
Per-task-type rates (\S\ref{sec:tasktypes}) are descriptive.

\subsection{Attribution of Bangla-only failures}
\label{sec:attribution}
For every EN-pass/BN-fail case we read both trajectories side by side. We record the \emph{first
failure point}, the earliest step that, corrected alone, would most likely have produced a pass,
and assign exactly one cause by checking, in order: (0) whether the Bangla output is already
correct for the task as written; (1) whether the translation is valid; (2) whether the simulated
user or coworker behaved correctly; (3--5) whether the agent understood the task and chose the
right tool and arguments; (6) whether the environment executed a valid action; and (7) whether the
agent verified its result and did not give up early (Figure~\ref{fig:overview}d). The first check
that fails decides the cause:

\begin{itemize}[noitemsep,topsep=1pt,leftmargin=*]
\item \textbf{Checker} (C5): the output is correct for the task as written but the grader rejects
it, or the check asks for something the task does not.
\item \textbf{Translation} (C3): the Bangla text changed the task or a literal.
\item \textbf{Simulated user} (C2): the Bangla simulated user or coworker broke the interaction.
\item \textbf{Tool / environment} (C4): a harness or tool fault blocked a valid action.
\item \textbf{Genuine Bangla failure} (C1): an agent error that the English run does not show.
\item \textbf{General weakness} (C6): an agent error whose mechanism also appears in an English
failure of the same model, cited by task.
\item \textbf{Unclear} (C7): no visible mechanism, or Bangla reruns of the task disagree.
\end{itemize}

\noindent \textbf{Checker--agent boundary.} A difference only in script or numerals (Bengali
digits, a Bangla word for a requested English keyword) is a checker artefact. A value changed into a
different English form (``curd'' for \emph{yogurt}, ``08:00 AM'' for \emph{8:00}) is an agent
failure. One \TAC{} case labelled translation (coworkers reply in Bangla and a checkpoint needs the
English word ``helpful'') is a checker artefact under this rule; moving it would change the
translation count from 9 to 8 and the checker count from 22 to 23, with the evaluation-side total
unchanged.

\paragraph{Reliability.} Labels were produced by annotation under a written codebook
(Appendix~\ref{app:codebook}). On \WB{}, a blind second annotator labelled every case and
disagreements were adjudicated (cause $\kappa=0.59$, $n=7$, \gemma{}; $\kappa=0.77$, $n=15$,
\gptoss{}). On \TAU{}, 46 of 66 cases were double-coded and adjudicated ($\kappa=0.39$ on the
fine cause, fair agreement on the scale of \citealp{landis1977}; $\kappa=0.71$, substantial, on
evaluation versus agent). On \OB{}, the codebook was checked on 16 double-coded cases
($\kappa=0.78$); the remaining cases and all \TAC{} cases have one annotator. The \WB{} and
\TAC{} annotators were LLMs \verify{annotator identity for \OB{} and \TAU{}; the draft's
$\kappa$ range 0.39--0.63 includes a value not found in the sources}.

\section{Results}

\subsection{Size of the Bangla Gap}
\label{sec:rq1}

\noindent\textbf{RQ1:} \emph{How large is the English--Bangla gap, and does it depend on the model and the benchmark?}

\noindent\textbf{Finding.} The Bangla gap depends on both the model and the benchmark. Only
\gemma{} loses significantly after correction, on \OB{} (8.0 points) and \TAU{} (11.5 points);
\gptoss{} is higher in Bangla on \OB{}, and neither model changes on \WB{}.

\begin{table*}[t]
\centering\small
\setlength{\tabcolsep}{6pt}
\begin{tabular}{llrrrcrrr}
\toprule
 & Model & $n$ & EN & BN & $\Delta$ [95\% CI] & MELR & $b$/$c$ & $p_{\mathrm{H}}$ \\
\midrule
OB  & gemma   & 300 & 50.3 & 42.3 & $+8.0$ [3.0, 13.0]        & 15.9   & 42/18 & \textbf{.016} \\
    & gpt-oss & 300 & 28.7 & 35.3 & $-6.7$ [$-$12.0, $-$1.3]  & $-$23.3 & 25/45 & .112 \\
WB  & gemma   & 90  & 81.1 & 82.2 & $-1.1$ [$-$10.0, 7.8]     & $-$1.4 & 7/8   & 1.00 \\
    & gpt-oss & 90  & 58.9 & 53.3 & $+5.6$ [$-$5.6, 16.7]     & 9.4    & 15/10 & 1.00 \\
$\tau^2$ & gemma & 218 & 49.5 & 38.1 & $+11.5$ [5.0, 17.4]   & 23.1   & 40/15 & \textbf{.007} \\
    & gpt-oss & 149 & 38.9 & 39.6 & $-0.7$ [$-$10.7, 8.7]     & $-$1.7 & 26/27 & 1.00 \\
TAC & gemma   & 23  & 39.1 & 26.1 & $+13.0$ [0.0, 26.1]       & 33.3   & 3/0   & 1.00 \\
\bottomrule
\end{tabular}
\caption{Success rates (\%) in English and Bangla on \OB{} (OB), \WB{} (WB), \TAU{} ($\tau^2$)
and \TAC{} (TAC). Positive $\Delta$ and MELR mean Bangla is worse. $b$/$c$: EN-only/BN-only tasks.
$p_{\mathrm{H}}$: exact McNemar test, Holm-corrected across the seven comparisons; bold:
$p_{\mathrm{H}}<.05$.}
\label{tab:rq1}
\end{table*}

Table~\ref{tab:rq1} gives every comparison. The two \gemma{} losses rest on 42 EN-only against
18 BN-only tasks on \OB{} ($50.3\%\to42.3\%$) and 40 against 15 on \TAU{} ($49.5\%\to38.1\%$).
\gptoss{}'s \OB{} gain ($28.7\%\to35.3\%$, 25 against 45 tasks) is significant before correction
($p=.022$) but not after. On \TAC{}, \gemma{} drops from 39.1\% to 26.1\%,
but this rests on three discordant tasks, so its MELR of 33.3\% is not informative; over the 21
unique tasks the gap is 9.5 points (2 against 0 tasks).

The two models also differ from each other. On \OB{}, \gemma{}'s gap exceeds \gptoss{}'s by
14.7 points (95\% CI [7.7, 22.3]). On the 149 \TAU{} tasks gradeable in all four cells, the gaps
are $+13.4$ and $-0.7$ points, a difference of 14.1 (95\% CI [2.0, 26.2]). The model gap itself
depends on the language. On \OB{} in English, \gemma{} solves 82 tasks that \gptoss{} fails and
\gptoss{} 17 that \gemma{} fails ($p<10^{-10}$); in Bangla the counts are 48 and 27 ($p=.020$).

\gptoss{}'s Bangla gain should not be read as better understanding. Its invalid actions fall
from 1{,}671 in English to 955 in Bangla over 299 runs each (Table~\ref{tab:trajectory}): the
Bangla prompt changes how well it follows \OB{}'s action format.

On \OB{}, \gemma{}'s gap is concentrated in short tasks: 16.1 points on single-app tasks (18
against 3 discordant tasks, Holm-adjusted $p=.004$), 9.5 on two-app tasks (16 against 7,
$p=.186$) and 0.0 on three-app tasks (8 against 8). Three-app tasks fail in both languages for
the same reasons, so they carry little language signal.

\subsection{Where Bangla Runs Break}
\label{sec:rq2}

\noindent\textbf{RQ2:} \emph{Where in the trajectory do Bangla-only failures first occur?}

\noindent\textbf{Finding.} Each benchmark breaks at a different point: Bangla-only failures start
at the final answer on \OB{}, in the dialogue on \TAU{}, at tool arguments on \WB{}, and at
verification and stopping on \TAC{}.

\begin{table}[t]
\centering\footnotesize
\setlength{\tabcolsep}{2.4pt}
\begin{tabular}{L{2.75cm}rrrrrrr}
\toprule
 & \multicolumn{2}{c}{OB} & \multicolumn{2}{c}{$\tau^2$} & \multicolumn{2}{c}{WB} & TAC \\
\cmidrule(lr){2-3}\cmidrule(lr){4-5}\cmidrule(lr){6-7}\cmidrule(lr){8-8}
First failure point & g & o & g & o & g & o & g \\
\midrule
Understanding, grounding, planning & 18 & 13 & 17 & 16 & 5 & 2  & 2 \\
Tool choice, arguments, execution & 6  & 3  & 10 & 3  & 2 & 12 & 2 \\
Verification, early stop & 3  & 6  & 12 & 4  & 0 & 1  & 3 \\
Final response only                     & 15 & 3  & 1  & 0  & 0 & 0  & 0 \\
Unclear                                 & 0  & 0  & 0  & 3  & 0 & 0  & 1 \\
\midrule
Cases                                   & 42 & 25 & 40 & 26 & 7 & 15 & 8 \\
\bottomrule
\end{tabular}
\caption{First failure point of the analysed cases (g: \gemma{}, o: \gptoss{}). The \TAC{}
cases include Bangla reruns that disagree and both-fail tasks with different partial scores.
On $\tau^2$, 10 of \gemma{}'s 11 early terminations were caused by the simulated customer.}
\label{tab:rq2}
\end{table}

Table~\ref{tab:rq2} gives the counts. On \OB{}, 15 of
\gemma{}'s 42 Bangla-only failures break only at the final response: the content is correct, but
it is written in Bangla script or digits, which the checker rejects. On \TAU{}, failures start
in grounding, planning and termination; 10 of \gemma{}'s 11 early terminations come from the
simulated customer ending the dialogue, not from the agent. On \WB{}, 11 of \gptoss{}'s 15
cases start at tool arguments: missing filters, wrong thresholds or invented argument names.
On \TAC{}, the analysed cases break at verification and stopping. No Bangla-only \TAU{} failure
starts with malformed tool syntax, and none of the 641 \WB{} tool-call arguments we inspected
contains Bangla script.

\subsection{Causes of Bangla-Only Failures}
\label{sec:rq3}

\noindent\textbf{RQ3:} \emph{Which pipeline component is responsible for each Bangla-only failure?}

\noindent\textbf{Finding.} The largest group of Bangla-only failures, 74 of 163 cases (45.4\%),
traces to evaluation components rather than the agent, and only 14.7--54.6\% can be
Bangla-specific agent failures. Correcting the grading cuts \gemma{}'s \OB{} gap from 8.0 to 3.0
points.

\begin{table}[t]
\centering\footnotesize
\setlength{\tabcolsep}{3pt}
\begin{tabular}{lrrrrr}
\toprule
Cause & OB & $\tau^2$ & WB & TAC & All (\%) \\
\midrule
\multicolumn{6}{l}{\emph{Evaluation side}} \\
Simulated user (C2) & 0 & 35 & 0 & 0 & 35 (21.5) \\
Checker (C5)        & 21 & 0 & 1 & 0 & 22 (13.5) \\
Translation (C3)    & 6 & 1 & 1 & 1 & 9 (5.5) \\
Tool / env. (C4)    & 7 & 0 & 0 & 1 & 8 (4.9) \\
\multicolumn{6}{l}{\emph{Agent side}} \\
Genuine Bangla (C1) & 13 & 8 & 3 & 0 & 24 (14.7) \\
General (C6)        & 4 & 21 & 17 & 3 & 45 (27.6) \\
Unclear (C7)        & 16 & 1 & 0 & 3 & 20 (12.3) \\
\midrule
Cases               & 67 & 66 & 22 & 8 & 163 \\
\bottomrule
\end{tabular}
\caption{Cause of every analysed EN-pass/BN-fail case, both models pooled per benchmark.
Evaluation-side causes account for 74 cases (45.4\%), agent causes for 69 (42.3\%) and unclear
cases for 20 (12.3\%).}
\label{tab:attribution}
\end{table}

Table~\ref{tab:attribution} attributes all 163 analysed cases. These are cases, not shares of the
gap: we did not apply the same codebook to every BN-pass/EN-fail task, so the table does not say
how much of each gap is artefactual. The evaluation-side cases come from the simulated user
(35, all on \TAU{}), checkers that reject correct answers (22,
almost all on \OB{}), translation (9) and tools or the environment (8). Agent causes account for
69 cases (42.3\%) and 20 are unclear. The Bangla-specific share is 14.7\% as labelled, 27.0\%
with the unclear cases and 54.6\% with the 45 general-weakness cases, whose label rests on one
English analogue per case that single runs cannot confirm.

The dominant cause follows the benchmark design: English-only checkers on \OB{}, the simulated
customer on \TAU{} and general weakness on \WB{}, where no simulated user is involved and the
grader compares database state. \TAC{} has too few cases for a dominant cause (3 general, 3
unclear, of which 2 are Bangla reruns that disagree).

\paragraph{Corrected grading and confound removal.} Table~\ref{tab:corrected} recomputes the gaps
under two corrections. Crediting answers that fail only on script or digits (on \OB{}, the 15
\gemma{} answers with correct content) cuts \gemma{}'s \OB{} gap from 8.0 to 3.0 points, and its
interval then includes zero. Removing evaluation-side cases moves \gemma{}'s \TAU{} gap from 11.5
to 5.8 points when dialogues that the simulated user ended before the agent acted are dropped in
\emph{both} languages. The corrections that touch only the Bangla side push the remaining gaps
down by construction, so we mark which ones are symmetric.

\begin{table*}[t]
\centering\small
\setlength{\tabcolsep}{8pt}
\begin{tabular}{llrrr}
\toprule
 & Model & As scored & Grading corr. & Confounds removed \\
\midrule
OB & gemma   & $+8.0$  & $+3.0$ [$-$1.3, 7.3]  & $+1.9$ [$-$1.9, 5.6]$^{s}$ \\
   & gpt-oss & $-6.7$  & $-8.7$ [$-$13.7, $-$3.3] & $-3.3$ [$-$7.7, 1.5]$^{s\dagger}$ \\
WB & gemma   & $-1.1$  & $-2.2$ [$-$10.1, 5.6]$^{s}$ & $-3.4$ [$-$11.4, 4.5]$^{s}$ \\
   & gpt-oss & $+5.6$  & $+6.7$ [$-$3.4, 18.0]$^{s}$ & $+8.0$ [$-$2.3, 18.2]$^{s}$ \\
$\tau^2$ & gemma & $+11.5$ & $+10.6$ [4.1, 17.0] & $+5.8$ [$-$1.9, 12.8]$^{s}$ \\
   & gpt-oss & $-0.7$  & $-0.7$ [$-$10.7, 8.7] & $-0.7$ [$-$10.5, 9.1]$^{s}$ \\
TAC & gemma  & $+13.0$ & $+8.7$ [0.0, 21.7]    & $0.0$ [0.0, 0.0]$^{s}$ \\
\bottomrule
\end{tabular}
\caption{Gap (points, 95\% CI) under corrected grading (OB, TAC: credit Bangla answers labelled
checker, C5; $\tau^2$: credit answers written with Bengali digits; WB: drop the ambiguous task) and after
removing evaluation-side cases. $^{s}$: applied to both languages (on OB, the English side covers
seven environment traps). $^{\dagger}$: also drops 19 BN-pass tasks whose English run failed on
action format.}
\label{tab:corrected}
\end{table*}

\paragraph{Two genuine Bangla mechanisms.} On \OB{}, all 13 genuine Bangla cases follow two
mechanisms, and both appear in both models. In \emph{entity grounding}, a name that appears only
in Bangla script has no letter case; the agent writes it in lower case for a case-sensitive mailbox
or calendar lookup, receives an empty result and accepts it (3 \gemma{} and 6 \gptoss{} cases;
the same two tasks fail for both models). In \emph{re-rendering}, the agent turns a value from
the Bangla task into a different English form: ``Trip'' for \emph{travelling} (both models),
``curd'' for \emph{yogurt}, and ``08:00 AM'' for \emph{8:00}. The same pattern appears once on
\WB{}, where ``last week'' is grounded as a rolling seven-day window. In each case the meaning
survives and the form the environment requires does not. The eight genuine \TAU{} cases are more
varied: grounding of names and identifiers, calls to user-side device tools in telecom, and policy
checks the English run performed.

\subsection{Quantified trajectory observations}
\label{sec:trajectory}

\begin{table}[t]
\centering\footnotesize
\setlength{\tabcolsep}{2.4pt}
\begin{tabular}{L{2.9cm}rrrr}
\toprule
 & \multicolumn{2}{c}{gemma} & \multicolumn{2}{c}{gpt-oss} \\
\cmidrule(lr){2-3}\cmidrule(lr){4-5}
Measure (\OB{}) & EN & BN & EN & BN \\
\midrule
Invalid actions (299 runs)        & 273 & 194 & 1{,}671 & 955 \\
Runs with $\geq$3 identical rejected actions in a row & 9 & 9 & 60 & 12 \\
Failed step repeated identically  & 11\% & 16\% & 17\% & 7\% \\
Writing runs that read back output & 3/180 & 4/178 & 20/112 & 19/155 \\
Out-of-step failures on 3-app tasks & 38/46 & 30/39 & 49/81 & 37/58 \\
\quad of which no step was rejected & 13/46 & 21/39 & 0/81 & 1/58 \\
\quad of which $\geq$half the steps were rejected & 9/46 & 10/39 & 44/81 & 9/58 \\
Bangla-script final answer        & -- & 46/173 & -- & 7/194 \\
\bottomrule
\end{tabular}
\caption{Trajectory measurements on \OB{}, computed from the run logs. Failed step repeated:
share of rejected or erroring steps followed by an identical action (284, 213, 1{,}640 and 949
failed steps). Read back: a written file is read again, or a created event is followed by
listing the calendar. Last row: share of Bangla failures.}
\label{tab:trajectory}
\end{table}

Table~\ref{tab:trajectory} reports patterns that we can count on \OB{}, the benchmark with full
logs in our repository; the stage counts below cover all failures, not only Bangla-only ones.

\paragraph{The two models fail at different stages.} An agent step moves from understanding to
planning, action encoding, execution, verification and completion. \gptoss{} often fails at
action encoding: 107 of its 214 English \OB{} failures produce no valid action or give up (54
coded as action format), against 12 of \gemma{}'s 149. On \WB{}, where function calling handles
the format, its failures move to the arguments: 40 of its 79 failed runs are wrong tool arguments,
against 11 of \gemma{}'s 33. \gemma{} usually reaches execution and fails later, on content (43 of
149 English failures), the step budget (42 of 149) and, in Bangla, the script of its answer.
The model gap follows the same line: of the 82 English tasks that \gemma{} solves and \gptoss{}
does not, 80 \gptoss{} runs contain a rejected action, 33 are coded action format, and 42 end
stuck or at the step limit.

\paragraph{Running out of steps means different things.} Out-of-step failures are mostly a
three-app problem (38 of 46 for \gemma{} in English). For \gemma{}, 13 of these 46 runs have no
rejected step at all: valid progress was too slow. For \gptoss{}, none of 81 is of this kind, and in
44 at least half of the steps were rejected, so the budget was spent on invalid actions.

\paragraph{Verification and repetition.} The agents rarely check what they wrote: \gemma{} reads
back its output in under 3\% of writing runs, \gptoss{} in 12--18\%. \gptoss{}'s English runs
contain 1{,}671 invalid actions and 60 runs with three or more identical rejected actions in a row;
both fall sharply in Bangla. One common claim is \emph{not} supported: most failed steps are not
followed by an identical action (7--17\% are), so repetition is concentrated in a minority of runs.
Whether missing verification \emph{causes} failures is a hypothesis (\S\ref{sec:discussion}).

\subsection{Evaluation-setup failures in both languages}
\label{sec:setup}

RQ3 counts only EN-pass/BN-fail cases. The evaluation setup also fails on tasks that fail in
both languages, which lowers both scores and can hide or mimic agent errors.

\paragraph{Scorer and benchmark defects (\OB{}).} The three scorer defects repaired in
\S\ref{sec:method} turned 19 English and 24 Bangla runs of \gemma{} (6 and 9 of \gptoss{}) from fail
to pass; the shipped scorer therefore understates both languages, Bangla slightly more. Beyond the
scorer, some checks ask for something the task does not: a benchmark defect is the primary cause of
47 of \gemma{}'s 149 English failures (31.5\%) and 39 of the 131 tasks that fail in both languages,
and it decides at least one of the three fully coded runs on 57 tasks. In task 1-2/0, all four runs
give the correct free time slot and all four fail, because two of the four checks require a
calendar event the task never asks for. Defects do not make a task uninformative: inside
defect-broken tasks the agents still make independent errors, so dropping these tasks would also
remove evidence about the agents.

\paragraph{Simulated users (\TAU{}).} A heuristic detector flags 46 \gemma{} Bangla dialogues and
26 English ones in which the simulated customer ended the chat before the agent acted (a hand check
of six random flags found about three true early stops). In 33 \gemma{} and 54 \gptoss{} Bangla
dialogues, the simulated customer wrote mostly in English, a fidelity problem the agent cannot fix.

\paragraph{Harness and environment limits.} On \OB{}, the in-container LLM helper answers in at most
50 tokens; for \gptoss{}, whose hidden reasoning uses that budget, all 14 helper calls in two Bangla
runs returned empty strings. The reply-length cap rejected 14 and 15 consecutive long writes in two
\gemma{} Bangla runs (Appendix~\ref{app:chains}), and two runs ended after empty model replies. On
\TAC{}, cloud-drive authentication errors and an administrator password prompt blocked valid actions
in both languages. On \TAU{}, \gptoss{} lost most retail and banking tasks to provider quota limits
and an unavailable judge model; these pairs are excluded on both sides (\S\ref{sec:method}). On
\WB{}, dropping every task with a provider error on either side leaves \gptoss{}'s gap almost
unchanged ($+7.1$ points, $n=85$).

\section{Analysis by Task Type}
\label{sec:tasktypes}

We assigned every task of \WB{}, \TAU{} and \TAC{} to one of nine types (Appendix~\ref{app:tasktypes},
Table~\ref{tab:percat}) \verify{who assigned types and how}. The results are descriptive: several
types have fewer than ten tasks, and we did not test per type. Three patterns hold for types with
at least ten tasks. First, failure rates differ more between types than between languages: on
\TAU{}, multi-target operations fail 73--93\% of the time in every cell ($n=44$ and 27). Second,
the largest increases are on \TAU{}, where the simulated customer speaks Bangla: single-target
tasks for \gemma{} ($35\%\to62\%$, $n=34$), entity creation for \gptoss{} ($55\%\to82\%$, $n=11$)
and conditional decisions for \gemma{} ($27\%\to43\%$, $n=37$). On \WB{}, \gemma{}'s failure
count differs by at most two tasks per type. Third, the two models can move in opposite directions
within a type: conditional decisions on \TAU{} change by $+16$ points for \gemma{} and $-10$ for
\gptoss{} ($n=37$ and 30), and system setup and troubleshooting by $+6$ and $-11$ ($n=87$ and
46). The opposite movement of the two models in \S\ref{sec:rq1} therefore persists inside task
types and is not only a product of task mix.

\section{Discussion}
\label{sec:discussion}

\paragraph{Implications for benchmark builders.} Our results suggest five practices.
(1) Separate checks of artefact state from checks of answer text, and accept the task language
or normalise numerals in answer checks; otherwise correct Bangla answers fail. (2) Constrain
simulated users: forbid confirming and ending a dialogue in the same turn, and require device
actions to be tool calls. (3) Protect names, literals and identifiers in translation, and audit
the result automatically. (4) Report discordant counts, corrected-grading gaps and per-case
attribution alongside raw scores. (5) Report results by benchmark and model; a pooled ``Bangla
penalty'' hides opposite effects.

\paragraph{Implications for agent builders.} The same failures point to scaffold changes that do
not touch the model: map common tool aliases to the benchmark's actions, re-prompt after a rejected
action instead of accepting a retry, allow bulk spreadsheet edits, require a read-back before
finishing, keep a short checklist of remaining sub-goals, and copy names, numbers and identifiers
verbatim across languages. We did not test these changes.

\paragraph{Hypotheses for future work.} The following explanations are consistent with our data
but not demonstrated by it. Each comes with the experiment that would test it.
\begin{itemize}[noitemsep,topsep=1pt,leftmargin=*]
\item \emph{Missing verification causes failures.} Test: require a read-back step before
finishing and measure the change in pass rate.
\item \emph{Long tasks fail by error accumulation.} Test: vary the step budget on three-app tasks
and compare with per-step error rates.
\item \emph{Scores measure model--scaffold compatibility.} \gptoss{}'s format failures suggest
it. Test: rerun with native function calling in both languages.
\item \emph{Script and tool-format errors share one cause}, a correct intent written in a form the
environment rejects. Test: a probe set with protected names and literals.
\item \emph{Early errors are rarely recovered.} Test: prompt the agent to retry after an empty
result and measure recovery.
\end{itemize}

\section{Conclusion}
We ran paired English and Bangla evaluations of two open-weight agents on four benchmarks. The
Bangla gap depends on model and benchmark: only two of seven comparisons are significant after
correction. Among the 163 tasks that fail only in Bangla, the largest group traces to evaluation
components, and correcting the grading removes most of one significant gap. The genuine failures
on \OB{} keep the meaning of the task but get its form wrong. Cross-lingual agent scores should be
read together with an attribution of their failures.

\section*{Limitations}

\paragraph{Case-level attribution.} We attribute cases, not gaps. On \OB{} and \TAU{} the
BN-pass/EN-fail tasks were not coded with the same codebook (\WB{} is the exception), so the
corrected gaps that drop cases only on the Bangla side are biased toward zero.

\paragraph{Single runs.} Each task ran once per language at temperature 0. Served inference is
not bit-deterministic: rerunning 13 Bangla \TAC{} tasks changed full completion for one and the
checkpoint score for three. The
general-weakness label assumes that one English analogue reflects a stable weakness.

\paragraph{Annotation.} Causes were assigned by one annotator per benchmark, with second review on
\WB{} (all cases), \TAU{} (46 of 66) and a 16-case check on \OB{}. Agreement on the fine cause is
fair to substantial. Some annotators were LLMs, which may share biases with the evaluated agents.

\paragraph{Translation.} Bangla tasks were machine-translated with LLM assistance. On \OB{}, the
translation rules were applied unevenly (\S\ref{sec:localization}), so some failures we attribute
to the agent may partly reflect translation. Native-speaker review of the translations is
incomplete.

\paragraph{Scope.} We study one language, two models and pilot-sized samples of \WB{} (90 tasks)
and \TAC{} (23 task pairs). On \TAU{} the two models have different task sets after infrastructure
exclusions, and the simulated user is played by the same model as the agent.

\section*{Ethical Considerations}
We evaluate public open-weight models on public benchmarks in sandboxed environments; no human
subjects or private data are involved. English-only graders can penalise an agent for answering in
the user's language, which understates agent quality for non-English users. We release task IDs,
trajectories, attribution labels and analysis scripts so that this bias can be measured. The
Bangla tasks were machine-translated and are not a curated Bangla language resource.

\bibliography{refs}

\appendix

\section{Attribution codebook}
\label{app:codebook}

\begin{table}[h]
\centering\footnotesize
\begin{tabular}{L{1.5cm}L{5.25cm}}
\toprule
Code & Operational definition \\
\midrule
C5 Checker & The Bangla output is correct for the task as written; the grader rejects it, or the check
requires something the task does not ask for. Includes script- or numeral-only differences. \\
C3 Translation & The Bangla task text differs from the English task in an operation, a constraint or a
literal (e.g., a quoted string translated; ``sum'' rendered as ``add''). \\
C2 Simulated user & The simulated user or coworker ends the interaction early, gives a wrong value,
replies in the wrong language, or narrates an action instead of performing it. \\
C4 Tool / env. & A tool, harness or infrastructure fault blocks an action that is valid (e.g., an
empty model reply recorded as malformed; a reply cap that rejects a long valid action). \\
C1 Genuine Bangla & The agent errs in a way the English run does not, and no check above fails
(e.g., a changed English value, a lower-cased name from Bangla script). \\
C6 General & The agent error's mechanism also appears in an English failure of the same model, cited
by task ID. \\
C7 Unclear & No visible mechanism, or Bangla reruns of the same task disagree. \\
\bottomrule
\end{tabular}
\caption{Causes, checked in the order C5, C3, C2, then agent understanding, tool choice and
arguments, C4, and verification (Figure~\ref{fig:overview}d). The first check that fails decides
the cause. Each case also records the first failure point, the decisive action and a trace quote.}
\label{tab:codebook}
\end{table}

\section{Reproducibility details}
\label{app:repro}

\begin{table}[h]
\centering\footnotesize
\begin{tabular}{L{1.8cm}L{4.75cm}}
\toprule
Item & Value \\
\midrule
Model tags & \texttt{gemma4:31b}, \texttt{gpt-oss:20b} \\
Serving & Ollama, OpenAI-compatible endpoint \verify{Ollama version, quantization, num\_ctx} \\
Decoding & temperature 0, top-$p$ 1 (\OB{} harness); \gptoss{} reasoning effort not set (default) \\
Reply cap (\OB{}) & 512 tokens (\gemma{}), 2{,}048 (\gptoss{}) \verify{per-model caps as stated in the \OB{} report} \\
Step budgets & \OB{} 20 steps; \TAC{} \verify{iteration cap}; \WB{}, \TAU{} \verify{limits} \\
Simulated users & \TAU{}: same model as the agent; \TAC{} coworkers \verify{model} \\
Translation & \OB{}: \gemma{}, cached; system prompt written in Bangla; others LLM-assisted \verify{tool and review} \\
\OB{} scorer & three repairs (\S\ref{sec:method}); scripts \texttt{analysis/fa\_rescore.py} \\
Statistics & exact McNemar, Holm, paired bootstrap (4{,}000 resamples, seed 7) \\
Annotation & codebook in Appendix~\ref{app:codebook}; LLM annotators on \WB{} and \TAC{} \verify{\OB{}, \TAU{}} \\
\bottomrule
\end{tabular}
\caption{Configuration. Fill in every value that the run logs and reports do not record before submission.}
\label{tab:repro}
\end{table}

\section{Task types}
\label{app:tasktypes}

\begin{table*}[t]
\centering\footnotesize
\setlength{\tabcolsep}{4pt}
\begin{tabular}{l rr rr rr rr rr}
\toprule
 & \multicolumn{2}{c}{WB gemma} & \multicolumn{2}{c}{WB gpt-oss} & \multicolumn{2}{c}{$\tau^2$ gemma}
 & \multicolumn{2}{c}{$\tau^2$ gpt-oss} & \multicolumn{2}{c}{TAC gemma} \\
\cmidrule(lr){2-3}\cmidrule(lr){4-5}\cmidrule(lr){6-7}\cmidrule(lr){8-9}\cmidrule(lr){10-11}
Task type & EN/BN & $n$ & EN/BN & $n$ & EN/BN & $n$ & EN/BN & $n$ & EN/BN & $n$ \\
\midrule
Entity creation          & 0/0   & 8  & 38/25 & 8  & 67/73 & 15 & 55/82 & 11 & 50/75   & 4 \\
Single-target            & 0/0   & 23 & 13/26 & 23 & 35/62 & 34 & 71/74 & 34 & 100/100 & 3 \\
Multi-target and batch   & 24/16 & 25 & 56/48 & 25 & 73/80 & 44 & 81/93 & 27 & -- & \\
Temporal reasoning       & 40/40 & 5  & 60/80 & 5  & -- & & -- & & 100/100 & 1 \\
Data analysis and QA     & 0/14  & 7  & 14/57 & 7  & 0/100 & 1  & 0/0   & 1  & 100/100 & 4 \\
Conditional decisions    & 42/42 & 19 & 63/68 & 19 & 27/43 & 37 & 37/27 & 30 & 100/100 & 2 \\
Cross-tool composition   & 33/33 & 3  & 33/33 & 3  & -- & & -- & & -- & \\
Interpersonal coord.     & -- & & -- & & -- & & -- & & 40/60 & 5 \\
System setup and troubl. & -- & & -- & & 53/59 & 87 & 61/50 & 46 & 0/25 & 4 \\
\midrule
All tasks                & 19/18 & 90 & 41/47 & 90 & 50/62 & 218 & 61/60 & 149 & 61/74 & 23 \\
\bottomrule
\end{tabular}
\caption{Failure rate (\%) in English/Bangla per task type, with the number of tasks. Descriptive:
with fewer than ten tasks one task moves a rate by 10--100 points.}
\label{tab:percat}
\end{table*}

\begin{table}[h]
\centering\footnotesize
\begin{tabular}{L{2.0cm}L{4.8cm}}
\toprule
Type & What the task asks \\
\midrule
Entity creation & Make one new record from details in the request \\
Single-target & Find the one record the request points to, then change it \\
Multi-target & Select every matching record, or make several changes \\
Temporal & Work out free slots, relative dates or time shifts \\
Data analysis & Compute, chart or look up an answer from data \\
Conditional & Check a condition or policy first; act only if it holds \\
Cross-tool & Use a value from one tool as input to another \\
Interpersonal & Obtain information or agreement from coworkers \\
System setup & Install software or walk a user through a fix \\
\bottomrule
\end{tabular}
\caption{The nine task types.}
\label{tab:types}
\end{table}

\section{Example failure chains}
\label{app:chains}

\begin{figure*}[h]
\centering
\includegraphics[width=\textwidth]{figures/fig_chains.pdf}
\caption{Four traced Bangla runs whose English counterpart passes, from the first divergence to
the final failure. These are illustrations; the counts are in Tables~\ref{tab:attribution}
and~\ref{tab:trajectory}.}
\label{fig:chains}
\end{figure*}

\end{document}
